\section{回顾：策略梯度的不足}

策略梯度使用采样的 reward-to-go 来估计每个动作的好坏，但这种估计\textbf{方差很大}。

\begin{warningbox}{策略梯度的两个核心问题}
\begin{enumerate}
    \item \textbf{数据利用不充分}：稀疏奖励下，大部分轨迹的奖励为零，梯度信息很少。
    \item \textbf{高方差}：采样 reward-to-go 是单次采样估计，噪声大。
\end{enumerate}
\end{warningbox}

核心改进思路：能否\textbf{学习}什么是好的、什么是坏的，而不是仅依赖采样估计？

\subsection{本章小结}

策略梯度的主要敌人是高方差与低效利用轨迹，构造一个学习的 critic 去估计优势函数可以显著改善这种现象，并为 Actor-Critic 系统的出现打下基础。

